\documentclass[french]{book}
\usepackage{teach}
\usepackage{amsmath,amsfonts,amssymb}
\usepackage{lscape}
%\usepackage[landscape]{geometry}
%\usepackage{geometry}
%\geometry{top=1cm, bottom=1cm, left=2cm, right=2cm}
\usepackage[utf8]{inputenc}

%\usepackage[latin1]{inputenc}
%\usepackage[T1]{fontenc}

\usepackage[french]{babel}
\redefinecolor{thm}{title}{bgcolor}{alizarin}
\redefinecolor{prop}{title}{bgcolor}{meatbrown}
\redefinecolor{defn}{title}{bgcolor}{olivine}
\definecolor{alizarin}{rgb}{0.82, 0.1, 0.26}
\definecolor{meatbrown}{rgb}{0.9, 0.72, 0.23}
\definecolor{olivine}{rgb}{0.6, 0.73, 0.45}
\usepackage{pgf,tikz}

\usepackage{mathrsfs}
\usetikzlibrary{arrows}
\usetikzlibrary{patterns}

\newcommand{\norme}[1]{\left\Vert #1\right\Vert}

\begin{document}
\pagestyle{empty}
\section*{Interprétation géométrique du déterminant}
%\chapter*{Interprétation géométrique du déterminant}
\begin{center}

\textit{L'objectif est de montrer dans un cas particulier que l'aire d'un parallélograme est égale}

 \textit{au déterminant des deux vecteurs formant le parallélogramme}
\end{center}
\subsection*{Construction}

Tracer les \underline{droites} $(OA);(OB);(AC);(BC)$.\\
Tracer un repère orthonormé $(O;\overrightarrow{i};\overrightarrow{j})$ où $\norme{\overrightarrow{i}}=\norme{\overrightarrow{j}}=1cm$ et tel que l'axe des abscisses soit parallèles à la feuille.\\

On note $\overrightarrow{u}=\overrightarrow{OB}$ et $\overrightarrow{v}=\overrightarrow{OA}$. On notera les coordonnées de $\overrightarrow{u}=\left( \begin{array}{c}
x_u\\
y_u
\end{array} \right)
$ et $\overrightarrow{v}=\left( \begin{array}{c}
x_v\\
y_v
\end{array} \right)$. \\

Tracer les droites $y=y_u; \; y=y_v; \; x=x_u; \; x=x_v$.\\

Placer : 
\begin{itemize}
\item W le point d'intersection de $y=y_v$ et l'axe des ordonnées
\item J le point d'intersection de $x=x_u$ et $(AC)$
\item P le point d'intersection de $(AC)$ et l'axe des ordonnées
\item K le point d'intersection de $y=y_v$ et $x=x_u$.\\
\end{itemize} 

À l'aide d'un compas, placer le point R tel que $\overrightarrow{OR}=\overrightarrow{AK}$. On note $R=(x_r;y_r)$ ces coordonnées.\\

Tracer la droite $x=x_r$.\\

Placer : 
\begin{itemize}
\item H le point d'intersection de $x=x_r$ et $OB$
\item T le point d'intersection de $x=x_r$  et $y=y_u$
\item Z le point d'intersection de $x=x_r$ et $y=y_v$.\\
\end{itemize} 

\subsection*{Interprétation}

Par construction, l'aire de AKJ est la même que ORH, l'aire de BCJ est la même que OAP, l'aire de HBT est la même que PAW. \\

Ainsi l'aire du parallélogramme OBCA est égale la somme de l'aire du rectangle ORZW et de l'aire du rectangle TBKZ.\\

\textit{
(Pour s'en persuader, vous pouvez photocopier votre figure et découper les triangles AKJ,BCJ et HBT puis les superposés respectivement sur les triangles ORH,OAP et PAW.)}\\


On sait que la longueur AK est égale à $x_u-x_v$. En déduire par contruction, la longueur de OR puis l'aire de ORZW sachant que la longueur OW est égale à $y_v$.\\

De même, on sait que la longueur TB est égale à $x_v$ et que la longueur BK est égale à $y_v-y_u$ en déduire l'aire de TBKZ.\\

Pour finir, déterminer l'aire du parallélogramme OBCA, puis conclure.
\newpage

\begin{landscape}
\begin{center}

\definecolor{ududff}{rgb}{0.30196078431372547,0.30196078431372547,1}
\begin{tikzpicture}[scale=1.5,line cap=round,line join=round,>=triangle 45,x=1cm,y=1cm]
\clip(-1.4915269060215481,-1.9491351810091735) rectangle (9.40069011843094,8.823551341117295);
\draw [->,line width=1pt] (0.39576749678791573,0.8506587593149608) -- (1.617831755373018,3.2549373550095515);
\draw [->,line width=1pt] (0.39576749678791573,0.8506587593149608) -- (6.26698926085982,3.0556877476315467);
\draw [->,line width=1pt] (1.617831755373018,3.2549373550095515) -- (7.489053519444922,5.459966343326137);
\draw [->,line width=1pt] (6.26698926085982,3.0556877476315467) -- (7.489053519444922,5.459966343326137);
\begin{scriptsize}
\draw [fill=ududff] (0.39576749678791573,0.8506587593149608) circle (1pt);
\draw[color=ududff] (0.09225224979733232,0.6282412753941756) node {$O$};
\draw [fill=ududff] (1.617831755373018,3.2549373550095515) circle (1pt);
\draw[color=ududff] (1.7357966567791894,3.6463864591244763) node {$A$};
\draw[color=black] (0.7795526381715634,2.1821378056315583) node {$\overrightarrow{v}$};
\draw [fill=ududff] (6.26698926085982,3.0556877476315467) circle (1pt);
\draw[color=ududff] (6.5170167498173175,2.824614255633553) node {$B$};
\draw[color=black] (3.4092236893425345,1.7534289242351887) node {$\overrightarrow{u}$};
\draw [fill=ududff] (7.489053519444922,5.459966343326137) circle (1pt);
\draw[color=ududff] (7.607732583541641,5.842759439363853) node {$C$};
\end{scriptsize}
\end{tikzpicture}

\end{center}
\end{landscape}
\end{document}